\documentclass[12pt,a4paper]{article}
\usepackage{amsmath,amsfonts,amssymb}
\usepackage{graphicx}
\usepackage{hyperref}
\usepackage{enumitem}
\usepackage{geometry}
\geometry{margin=2.5cm}
\emergencystretch=2em

\begin{document}

\begin{center}
    {\Large \textbf{Master Thesis Project in Quantum Information Science and Technology}}\\[0.3cm]
    {\large \textbf{Parametric Matrix Models for Quantum Dynamics: Trotterized Time Evolution, Unitary Coupled Cluster Circuits and Time-Dependent Hartree--Fock}}\\[0.2cm]
    \vspace{0.3cm}
    {Duration: One Year (60 ECTS)}
\end{center}

\vspace{1cm}

\section*{1. Introduction}

\emph{Parametric matrix models} (PMMs)~[1] are a class of
machine-learning models built from matrix equations instead of
neurons.  The trainable parameters are the elements of Hermitian or
unitary matrices that depend on the input features, and the outputs
are eigenvalues of these matrices or expectation values in their
eigenvectors.  A PMM therefore has the same mathematical form as the
problems of quantum mechanics.  It can be seen as eigenvector
continuation~[2], a reduced-basis method in which a large Hamiltonian
is projected on a few of its eigenvectors, with the difference that
the elements of the small matrices are learned directly from data.
PMMs are universal function approximators, they need few parameters,
their outputs have the analytic structure of eigenvalue problems, and
in the applications of Ref.~[1] they extrapolated well outside the
training data.  One of these applications was the extrapolation of the
energies of a Trotterized spin chain to zero time step.

Many problems in quantum many-body physics and quantum computing
consist of a unitary evolution that is expensive to compute or to run,
and that has to be repeated for many values of a parameter.  This
thesis project studies PMMs as models of such evolutions.  Three
settings are considered, and they define three possible paths.
\begin{itemize}
  \item[\textbf{A.}] \textbf{Trotterized time evolution.}  On a
  quantum computer the evolution operator $e^{-iHt}$ is replaced by a
  product of simpler exponentials with a finite time step~[3,4].  A
  smaller step gives a smaller error but a deeper circuit and more
  noise.  A PMM with the same product structure is trained on data
  taken at moderate time steps, from simulations and from quantum
  hardware, and is used to extrapolate to zero time step.
  \item[\textbf{B.}] \textbf{Trotterized unitary coupled cluster.}
  The unitary coupled cluster (UCC) ansatz~[5] for the pairing
  Hamiltonian~[6,7] is implemented as a product of exponentials of
  pair excitations.  A PMM is trained to reproduce the energies and
  observables of these circuits as functions of the number of Trotter
  steps, the pairing strength and the amplitudes, and to extrapolate
  them.
  \item[\textbf{C.}] \textbf{Time-dependent Hartree--Fock.}  The
  time-dependent Hartree--Fock (TDHF) equations~[8] describe the
  mean-field dynamics of nuclei, atoms and electrons, and are costly
  when they have to be solved for many interaction parameters and
  initial conditions.  A PMM is constructed as a reduced model of
  these nonlinear matrix equations.
\end{itemize}
The three paths share the PMM theory and one PMM code, which are
developed in a common first part together with a reproduction of
results from Ref.~[1].  A thesis is expected to carry out the common
part and one path in depth, possibly with elements of a second.
Paths~A and~B use the same Trotter formalism and combine naturally.
Path~A includes experiments on real quantum computers, Path~B is
closest to nuclear many-body theory, and Path~C is a project in
classical computational physics.  The choice is made together with the
supervisors after the first months.

\section*{2. Parametric Matrix Models and Scientific Motivation}

\textbf{Parametric matrix models.}  A PMM consists of primary
matrices, which are Hermitian or unitary and whose elements are
analytic functions of the input features
$\boldsymbol c=(c_1,c_2,\dots)$, and of secondary matrices $S_k$.  The
outputs are eigenvalues of the primary matrices and bilinear forms
$\boldsymbol v^\dagger S_k\boldsymbol v$ of their eigenvectors
$\boldsymbol v$.  The simplest case is the affine eigenvalue PMM,
\[
  M(\boldsymbol c)=M_0+\sum_l c_l\,M_l ,
\]
with trainable Hermitian $n\times n$ matrices $M_l$ and the
eigenvalues $E_k(\boldsymbol c)$ as outputs.  This is the form of a
Hamiltonian that depends linearly on its coupling constants, reduced
to dimension $n$.  Its eigenvalues have branch points in the complex
plane of the features and avoided crossings on the real axis, which a
network of standard activation functions has to build from many
neurons.  For unitary evolution the primary matrix is a product of
exponentials,
\[
  U(dt)=\prod_{l=1}^{L}\exp\bigl(-iM_l\,dt\bigr),
\]
with eigenvalues $e^{-iE_k(dt)\,dt}$.  For $dt\to0$ the energies
$E_k$ become the eigenvalues of $\sum_lM_l$, so that a model trained
at finite time steps gives the result at zero time step directly.  In
Ref.~[1] five factors of dimension $9\times9$ reproduced the lowest
energies of a chain of ten spins with Heisenberg and
Dzyaloshinskii--Moriya interactions, and the extrapolation to zero
step was more accurate than polynomial extrapolation and than a neural
network.  For time-dependent observables a natural form is
\[
  \langle O\rangle(t;\boldsymbol c)\simeq
  \boldsymbol\phi_0^\dagger\,U^\dagger(t;\boldsymbol c)\,S\,
  U(t;\boldsymbol c)\,\boldsymbol\phi_0 ,
\]
with a trainable initial vector $\boldsymbol\phi_0$, a secondary
matrix $S$ for each observable and a propagator $U$ built from
primary matrices as above.  Unitarity, and with it the conservation
of the norm, is exact by construction.

\textbf{Path A: Trotter errors and effective Hamiltonians.}  For
$H=\sum_jH_j$ the first-order product formula defines an effective
Hamiltonian,
\[
  \prod_j e^{-iH_j\,dt}=e^{-iH_{\rm eff}(dt)\,dt},\qquad
  H_{\rm eff}(dt)=H-\frac{i\,dt}{2}\sum_{j<k}[H_j,H_k]+O(dt^2),
\]
and higher-order formulas remove the leading terms at the cost of more
gates.  Energies measured with phase estimation or the rodeo
algorithm, and observables after $n=t/dt$ steps, are those of
$H_{\rm eff}(dt)$.  Standard remedies are extrapolations in $dt$ with
polynomials~[9].  A PMM with the product structure of the circuit
learns small matrices that play the role of the terms $H_j$, and
thereby a model of $H_{\rm eff}(dt)$ for all $dt$.  On hardware the
gain is that the data are taken where the circuits are shallow.  The
open issues are how the required matrix dimension and number of data
points grow with the size and the entanglement of the system, and
whether the extrapolation survives the noise of a real device.

\textbf{Path B: Unitary coupled cluster for the pairing Hamiltonian.}
The pairing Hamiltonian,
\[
  H=\sum_{p}\varepsilon_p N_p-G\sum_{p,q}P_p^\dagger P_q ,
\]
with single-particle energies $\varepsilon_p$, number operators $N_p$,
pair creation operators $P_p^\dagger$ and pairing strength $G$, is a
minimal model of nuclear pairing and superconductivity.  Without
broken pairs it maps to one qubit per level, and it is exactly
solvable~[6].  With a reference state $|\Phi_0\rangle$ in which the
lowest levels are occupied by pairs, the UCC ansatz with pair
excitations is
\[
  |\Psi(\boldsymbol\theta)\rangle
  =\exp\Bigl(\sum_j\theta_j\tau_j\Bigr)|\Phi_0\rangle,\qquad
  \tau_j=P_a^\dagger P_i-P_i^\dagger P_a ,
\]
where $j=(i,a)$ runs over occupied levels $i$ and empty levels $a$.
Since the $\tau_j$ do not commute, the circuit implements
\[
  U_n(\boldsymbol\theta)
  =\Bigl[\prod_j\exp\bigl(\theta_j\tau_j/n\bigr)\Bigr]^{n},
\]
and the energy
$E_n(\boldsymbol\theta;G)=\langle\Phi_0|U_n^\dagger HU_n|\Phi_0\rangle$
depends on the number of Trotter steps $n$.  The limit $n\to\infty$ is
the UCC ansatz, while each finite $n$ is a variational ansatz of its
own when the amplitudes are optimized again~[5].  The PMM is to learn
the dependence on $1/n$, $G$ and $\boldsymbol\theta$ from circuits
with few steps and to extrapolate in $1/n$ and in $G$, including the
region where the system changes from a normal to a paired state.

\textbf{Path C: Time-dependent Hartree--Fock.}  The TDHF equation for
the one-body density matrix,
\[
  i\hbar\,\frac{d\rho}{dt}=\bigl[h[\rho(t)],\rho(t)\bigr],\qquad
  \rho(t+dt)=U(dt)\,\rho(t)\,U^\dagger(dt),\quad
  U(dt)=\exp\Bigl(-\frac{i}{\hbar}\,h[\rho(t)]\,dt\Bigr),
\]
is a unitary evolution with a single-particle Hamiltonian that depends
on the density itself.  It is the basis for the description of
collective motion, giant resonances, fusion and fission in nuclei, and
of mean-field dynamics in quantum chemistry and cold atoms~[8].  Each
time step requires the construction of $h[\rho]$ and a matrix
exponential, and studies over interaction strengths, trap parameters
and initial conditions require many complete simulations.  The first
test system is the Lipkin--Meshkov--Glick model~[10],
\[
  H=\varepsilon J_z+\tfrac12V\bigl(J_+^2+J_-^2\bigr),
\]
for which TDHF reduces to the motion of a classical spin, the
Hartree--Fock state changes character at $\chi=(N-1)|V|/\varepsilon=1$,
and exact solutions are available for comparison.  The point of
principle is that TDHF is nonlinear, while a PMM with fixed primary
matrices gives a linear evolution in its latent space.  Two
constructions will be compared: a linear unitary PMM in a latent space
larger than the single-particle space, and a self-consistent PMM whose
primary matrix depends on the latent density $\tilde\rho$,
\[
  M[\tilde\rho]=M_0+\sum_l\mathrm{Tr}\bigl(\tilde\rho\,S_l\bigr)\,M_l ,
\]
which has the structure of the Hartree--Fock Hamiltonian.

\textbf{Common questions.}  The methods are compared through the
accuracy of the extrapolation (in the time step, in the parameters and
in time), the amount of training data, the size of the model (matrix
dimension and number of parameters), and the conservation of energy
and particle number.  The project is organized around four questions:
\begin{enumerate}
  \item How small can a PMM be, and how little data does it need, for
  a reliable extrapolation to zero Trotter step, compared with
  polynomial extrapolation and with neural networks?  How does this
  change with system size, entanglement and the order of the product
  formula?
  \item Does the extrapolation work with data from noisy quantum
  hardware, and can it be combined with error mitigation?
  \item Can a PMM reproduce Trotterized UCC energies and observables
  of the pairing Hamiltonian over the pairing strength and the
  amplitudes, and what does this save in circuit depth?
  \item Can a small PMM reproduce nonlinear mean-field trajectories
  over a range of parameters and initial conditions, and how do its
  errors grow in time?
\end{enumerate}
In all cases the learned matrices are to be examined as well: whether
their spectra resemble the low-lying spectra of the physical
Hamiltonians, and what their analytic structure is.

\section*{3. Objectives of the Thesis}

\subsection*{Objective 1: PMM theory, code and reproduction (common part)}
\begin{itemize}
  \item Study the PMM framework of Ref.~[1]: Hermitian and unitary
  primary matrices, eigenvalue-based and observable-based models,
  universal approximation, and the relation to eigenvector
  continuation and reduced-basis methods~[2].
  \item Implement a PMM code with affine eigenvalue models, observable
  models and unitary product models, trained with the analytic
  gradients of Ref.~[1] or with automatic differentiation for complex
  matrices.
  \item Reproduce results of Ref.~[1]: the emulation of energies and
  observables of a model Hamiltonian as functions of a coupling, and
  the extrapolation to zero Trotter step for the spin chain, with
  polynomial extrapolation and a multilayer perceptron as baselines.
  \item Study Trotter and Suzuki--Trotter formulas and their
  errors~[3,4].
\end{itemize}

\subsection*{Objective 2 (Path A): PMMs for Trotterized quantum dynamics}
\begin{itemize}
  \item Generate Trotterized data (energies of the product formula and
  time-dependent observables) by classical simulation for the
  transverse-field Ising model, Heisenberg and XXZ chains, the
  Fermi--Hubbard model with two to four qubits, and the
  Lipkin--Meshkov--Glick model.
  \item Construct PMMs with the product structure of first- and
  second-order formulas, train them at moderate $dt$, extrapolate to
  $dt\to0$ and extract the effective Hamiltonian $H_{\rm eff}(dt)$.
  \item Compare with exact results, with polynomial extrapolations of
  the Trotter error~[9] and with neural networks, as functions of
  system size, entanglement, step size and Hamiltonian.
  \item Execute Trotterized circuits for selected Hamiltonians on real
  quantum hardware for a range of $dt$, train PMMs on the measured
  data, and determine whether they extrapolate to $dt\to0$, learn
  $H_{\rm eff}(dt)$ and compensate for hardware noise.
  \item Combine the extrapolation in $dt$ with error mitigation~[11],
  for example with a second input feature for the noise level.
\end{itemize}

\subsection*{Objective 3 (Path B): PMMs for Trotterized unitary coupled cluster}
\begin{itemize}
  \item Implement the pairing Hamiltonian in its qubit form and its
  exact solution for a few levels and pairs, and review UCC and
  pair-specific variants of it.
  \item Build and simulate the Trotterized UCC circuits, and generate
  energies and observables (pair occupations, pairing gap) as
  functions of the number of Trotter steps, the pairing strength $G$,
  the level spacing and the amplitudes.
  \item Design PMMs for these data: primary matrices that are affine
  in $1/n$, $G$ and the amplitudes, or that have the product structure
  of the circuit, with secondary matrices for the observables.
  \item Train on circuits with few Trotter steps and extrapolate to
  $n\to\infty$; train on a range of $G$ and extrapolate from weak to
  strong pairing.  Compare with polynomial extrapolation, splines and
  multilayer perceptrons.
  \item Determine the cost: number of parameters, number of circuit
  evaluations needed for training, and the circuit depth saved.
\end{itemize}

\subsection*{Objective 4 (Path C): PMMs for time-dependent Hartree--Fock}
\begin{itemize}
  \item Implement a TDHF solver (Runge--Kutta, Crank--Nicolson, Magnus
  expansion or a split-operator method), test its convergence in the
  time step and the basis size, and generate trajectories of density
  matrices and observables.
  \item Start with the Lipkin--Meshkov--Glick model and continue with
  two-level pairing models, a harmonic-oscillator trap with mean-field
  interactions, and small Hubbard or spin-chain systems.
  \item Construct the linear and the self-consistent PMM of Section~2,
  with primary matrices for the effective Hamiltonian or the
  propagator and secondary matrices for the observables, and study
  how commutation relations of the many-body Hamiltonian can be
  imposed as constraints.
  \item Train the models to reproduce trajectories as functions of
  interaction strengths, model parameters and initial conditions, and
  study overfitting, stability and extrapolation in the parameters and
  in time.
  \item Compare with the full solver, with recurrent neural networks,
  with dynamic mode decomposition, with reduced-basis methods and, in
  collaboration with other students, with physics-informed neural
  networks~[12].  Evaluate the growth of the error in time and the
  conservation of norm, energy and particle number.
\end{itemize}

\subsection*{Objective 5: Performance analysis and interpretation (all paths)}
\begin{itemize}
  \item Quantify the efficiency of the PMMs: number of trainable
  parameters, data requirements, and cost of training and evaluation,
  against the baselines of each path.
  \item Analyze what the models have learned: spectra of the primary
  matrices compared with the physical spectra, avoided crossings as
  functions of $dt$, $G$ or the interaction strength, and the
  analytic continuation to complex values of the features.
  \item Establish where the models fail: time steps beyond the radius
  of convergence, long times, strong entanglement, and parameters far
  from the training region.
\end{itemize}

\subsection*{Objective 6: Possible extensions}
\begin{itemize}
  \item PMMs for error mitigation, with noise-adapted effective
  Hamiltonians, and PMMs as surrogates in hybrid classical--quantum
  workflows.
  \item Exceptional points and analytic continuation of the learned
  models.
  \item Trotterized time evolution inside phase estimation or the
  rodeo algorithm, with a PMM extrapolation of the measured energies
  to zero time step.
  \item Open quantum systems, in connection with the tensor PMM for
  dissipative dynamics developed in the group~[13].
  \item PMM emulation of TDHF linear-response spectra and of TDHF in
  higher dimensions; PMMs trained on exact dynamics beyond the mean
  field.
  \item PMM surrogates for the time-dependent simulations of electrons
  in multi-well traps~[14], as fast environments for optimal control
  and reinforcement learning.
\end{itemize}

\section*{4. Methods and Tools}

\begin{itemize}
  \item \textbf{PMM code:} Python with NumPy and SciPy, and PyTorch or
  JAX for automatic differentiation; the examples of Ref.~[1] serve
  as reference.
  \item \textbf{Quantum circuits (Paths~A and~B):} Qiskit or
  PennyLane with state-vector and noisy simulators; IBM Quantum, AWS
  Braket or Azure Quantum for the hardware runs of Path~A.
  \item \textbf{TDHF (Path~C):} Python for the model systems, with
  C++ or Fortran for the time-critical parts of larger systems.
  \item \textbf{Useful background:} quantum mechanics, many-body
  physics (FYS4480), machine learning (FYS-STK4155) and quantum
  computing (FYS5419), with the lecture notes and codes of
  Ref.~[15].
\end{itemize}
All codes and results will be documented in Jupyter notebooks so that
the thesis can be reproduced.

\section*{5. Work Plan and Milestones (One-Year Timeline)}

\subsection*{Semester 1 (Months 1--6)}

\textbf{Months 1--2: Theory and background}
\begin{itemize}
    \item PMM theory, Trotter formulas, and the theory of the chosen
    path (product formulas and hardware, pairing and UCC, or TDHF).
\end{itemize}
\textbf{Deliverable:} Written summary of the theoretical background.

\bigskip

\textbf{Months 3--4: PMM code and reproduction}
\begin{itemize}
    \item PMM code with Hermitian and unitary models; reproduction of
    the examples of Ref.~[1]; tests on Hamiltonians with one and two
    qubits; choice of path.
\end{itemize}
\textbf{Deliverable:} Verified PMM code and reproduced reference
results.

\bigskip

\textbf{Months 5--6: Data generation and first models}
\begin{itemize}
    \item Path~A: Trotterized simulations of the Ising, XXZ and
    Lipkin--Meshkov--Glick models; first extrapolations.
    \item Path~B: pairing Hamiltonian and Trotterized UCC circuits;
    data for energies and observables against the number of steps.
    \item Path~C: TDHF solver and trajectories for the
    Lipkin--Meshkov--Glick model; first PMMs with small matrices on
    short trajectories.
\end{itemize}
\textbf{Deliverable:} Data sets and a first working PMM for the chosen
path.

\subsection*{Semester 2 (Months 7--12)}

\textbf{Months 7--9: Full training and extrapolation studies}
\begin{itemize}
    \item Path~A: experiments on quantum hardware; PMMs trained on
    device data; extrapolation to zero step.
    \item Path~B: data over pairing strengths and amplitudes;
    extrapolation in the number of steps and in $G$.
    \item Path~C: training over parameter domains; long-time stability
    and conservation laws; further model systems.
\end{itemize}
\textbf{Deliverable:} Trained models and a numerical analysis of the
extrapolation quality.

\bigskip

\textbf{Months 10--11: Benchmarking and interpretation}
\begin{itemize}
    \item Comparison with the baselines; efficiency and robustness;
    spectra and analytic structure of the learned matrices; one of the
    extensions of Objective~6, depending on progress.
\end{itemize}
\textbf{Deliverable:} Benchmark results and drafts of the results
chapters.

\bigskip

\textbf{Month 12: Writing and final submission}
\begin{itemize}
    \item Completion of the thesis manuscript (methods and theory
    chapters are written continuously during the year) and
    preparation of the final presentation.
\end{itemize}
\textbf{Deliverable:} Completed thesis and defense presentation.

\section*{6. Expected Outcomes}

\begin{itemize}
    \item A documented PMM code with Hermitian, unitary and
    time-evolution models.
    \item For the chosen path: PMMs for the extrapolation of
    Trotterized dynamics to zero time step, including a demonstration
    with data from quantum hardware; PMMs for Trotterized UCC
    simulations of the pairing Hamiltonian; or a PMM surrogate for
    nonlinear mean-field dynamics, together with a TDHF code and data
    set.
    \item Quantitative benchmarks of PMMs against polynomial
    extrapolation, neural networks and reduced-basis methods, in
    accuracy, data requirements and number of parameters.
    \item Insight into what the learned matrices represent, and into
    how physical structure can be built into machine-learning models.
    \item A master thesis and, depending on the results, a
    contribution to a scientific article.
\end{itemize}

\section*{7. References}

\begin{enumerate}
    \item P.~Cook, D.~Jammooa, M.~Hjorth-Jensen, D.~D.~Lee and D.~Lee, \textit{Parametric matrix models}, Nat.\ Commun.\ \textbf{16}, 5929 (2025).
    \item D.~Frame, R.~He, I.~Ipsen, D.~Lee, D.~Lee and E.~Rrapaj, \textit{Eigenvector continuation with subspace learning}, Phys.\ Rev.\ Lett.\ \textbf{121}, 032501 (2018); S.~K\"onig, A.~Ekstr\"om, K.~Hebeler, D.~Lee and A.~Schwenk, \textit{Eigenvector continuation as an efficient and accurate emulator for uncertainty quantification}, Phys.\ Lett.\ B \textbf{810}, 135814 (2020); T.~Duguet, A.~Ekstr\"om, R.~J.~Furnstahl, S.~K\"onig and D.~Lee, \textit{Colloquium: Eigenvector continuation and projection-based emulators}, Rev.\ Mod.\ Phys.\ \textbf{96}, 031002 (2024).
    \item H.~F.~Trotter, \textit{On the product of semi-groups of operators}, Proc.\ Amer.\ Math.\ Soc.\ \textbf{10}, 545 (1959); M.~Suzuki, \textit{Generalized Trotter's formula and systematic approximants of exponential operators and inner derivations with applications to many-body problems}, Commun.\ Math.\ Phys.\ \textbf{51}, 183 (1976).
    \item A.~M.~Childs, Y.~Su, M.~C.~Tran, N.~Wiebe and S.~Zhu, \textit{Theory of Trotter Error with Commutator Scaling}, Phys.\ Rev.\ X \textbf{11}, 011020 (2021).
    \item A.~Anand \textit{et al.}, \textit{A quantum computing view on unitary coupled cluster theory}, Chem.\ Soc.\ Rev.\ \textbf{51}, 1659 (2022); H.~R.~Grimsley, D.~Claudino, S.~E.~Economou, E.~Barnes and N.~J.~Mayhall, \textit{Is the Trotterized UCCSD Ansatz Chemically Well-Defined?}, J.\ Chem.\ Theory Comput.\ \textbf{16}, 1 (2020).
    \item R.~W.~Richardson, \textit{A restricted class of exact eigenstates of the pairing-force Hamiltonian}, Phys.\ Lett.\ \textbf{3}, 277 (1963); J.~Dukelsky, S.~Pittel and G.~Sierra, \textit{Colloquium: Exactly solvable Richardson--Gaudin models for many-body quantum systems}, Rev.\ Mod.\ Phys.\ \textbf{76}, 643 (2004).
    \item D.~J.~Dean and M.~Hjorth-Jensen, \textit{Pairing in nuclear systems: from neutron stars to finite nuclei}, Rev.\ Mod.\ Phys.\ \textbf{75}, 607 (2003).
    \item P.~Ring and P.~Schuck, \textit{The Nuclear Many-Body Problem}, Springer (1980); C.~Simenel, \textit{Nuclear quantum many-body dynamics: From collective vibrations to heavy-ion collisions}, Eur.\ Phys.\ J.\ A \textbf{48}, 152 (2012); T.~Nakatsukasa, K.~Matsuyanagi, M.~Matsuo and K.~Yabana, \textit{Time-dependent density-functional description of nuclear dynamics}, Rev.\ Mod.\ Phys.\ \textbf{88}, 045004 (2016).
    \item S.~Endo, Q.~Zhao, Y.~Li, S.~Benjamin and X.~Yuan, \textit{Mitigating algorithmic errors in a Hamiltonian simulation}, Phys.\ Rev.\ A \textbf{99}, 012334 (2019); G.~Rendon, J.~Watkins and N.~Wiebe, \textit{Improved Accuracy for Trotter Simulations Using Chebyshev Interpolation}, Quantum \textbf{8}, 1266 (2024); Yang \textit{et al.}, Phys.\ Rev.\ A \textbf{106}, 042401 (2022).
    \item H.~J.~Lipkin, N.~Meshkov and A.~J.~Glick, \textit{Validity of many-body approximation methods for a solvable model (I)}, Nucl.\ Phys.\ \textbf{62}, 188 (1965).
    \item K.~Temme, S.~Bravyi and J.~M.~Gambetta, \textit{Error Mitigation for Short-Depth Quantum Circuits}, Phys.\ Rev.\ Lett.\ \textbf{119}, 180509 (2017); Z.~Cai \textit{et al.}, \textit{Quantum error mitigation}, Rev.\ Mod.\ Phys.\ \textbf{95}, 045005 (2023).
    \item S.~Hochreiter and J.~Schmidhuber, \textit{Long Short-Term Memory}, Neural Comput.\ \textbf{9}, 1735 (1997); P.~J.~Schmid, \textit{Dynamic mode decomposition of numerical and experimental data}, J.\ Fluid Mech.\ \textbf{656}, 5 (2010); M.~Raissi, P.~Perdikaris and G.~E.~Karniadakis, \textit{Physics-informed neural networks}, J.\ Comput.\ Phys.\ \textbf{378}, 686 (2019).
    \item J.~D.~\'Alvarez-Cuartas, J.~H.~Reina, D.~Jammooa, M.~Hjorth-Jensen and D.~Lee, \textit{Tensor Parametric Matrix Model and the Dissipative Many-Body Problem}, in preparation (2026).
    \item O.~Leinonen \textit{et al.}, \textit{Design and dynamics of two-qubit gates with motional states of electrons on helium}, Phys.\ Rev.\ A \textbf{113}, 032624 (2026).
    \item M.~Hjorth-Jensen, \textit{Quantum Computing and Quantum Machine Learning}, lecture notes and codes,\\ \url{https://github.com/CompPhysics/QuantumComputingMachineLearning}.
\end{enumerate}

\end{document}
